\chapter{Programs — Cells}
%%% generated by the blueprint pipeline from TCSlib.Programs.Light.Compiler.Cells
\section{Overview}

Declarations of the Lean module \texttt{TCSlib.Programs.Light.Compiler.Cells}.

\section{Declarations}

\begin{definition}[Word size fits the limits]
\lean{Light.Compiler.Fits}
\label{Light.Compiler.Fits}
\leanok
\uses{Light.Limits}
Let $W,N,Q\in\mathbb N$ and let the limits have word bound $w\in\mathbb Z$ and space
bound $s\in\mathbb N$. The word size fits these limits, pool size $N$, and stack size
$Q$ when $2(2w+2)<2^W$, $2s<2^W$, $4Q<2^W$, and $2N<2^W$, with the natural quantities
interpreted as integers in these inequalities.
\end{definition}

\begin{theorem}[Word values lie in range]
\lean{Light.Compiler.Fits.inRange}
\label{Light.Compiler.Fits.inRange}
\leanok
\uses{Light.Compiler.Fits, Light.Limits, WordRAM.InRange, WordRAM.inRange_of_abs_le}
\difficulty{3}
Let $W,N,Q$ be natural numbers, let the limits have word bound $w$, and suppose
$2(2w+2)<2^W$, $2\,\mathrm{space}<2^W$, $4Q<2^W$, and $2N<2^W$. If an integer $v$
satisfies $|v|\le w$, then $-2^W\le 2v<2^W$.
\end{theorem}

\begin{definition}[Compiler size configuration]
\lean{Light.Compiler.Sizes}
\label{Light.Compiler.Sizes}
\leanok
\uses{Light.Compiler.CodeLayout, Light.Compiler.Fits, Light.Limits}
For a natural-number word size $W$, a size configuration consists of a code layout with
$F$ local variables per frame, a list of procedure entry positions, and a dispatcher
position $D$, together with execution limits and a natural-number stack size $Q$. The
word size must satisfy the compatibility condition relating $W$, the execution limits,
$D$, and $Q$, and the dispatcher position must satisfy $2(F+1)\le D$.
\end{definition}

\begin{definition}[Temporary addresses]
\lean{Light.Compiler.Temps}
\label{Light.Compiler.Temps}
\leanok
\uses{Light.Compiler.cADDR, Light.Compiler.cT}
For natural numbers $k$ and $F$, the temporary-address set consists of the distinguished
address $c_{\mathrm{ADDR}}$ together with every address $c_T(j)$ for which $k\le j\le
2F$.
\end{definition}

\begin{definition}[Stack addresses]
\lean{Light.Compiler.Stack}
\label{Light.Compiler.Stack}
\leanok
\uses{Light.Compiler.cStack}
For natural numbers $q$ and $Q$, the stack-address set consists of the addresses
$c_{\mathrm{Stack}}(j)$ with $q\le j<Q$.
\end{definition}

\begin{definition}[Main frame address set]
\lean{Light.Compiler.MainFrame}
\label{Light.Compiler.MainFrame}
\leanok
\uses{Light.Compiler.Stack, Light.Compiler.mainFrame}
For a natural number $F$, let the main frame address set consist of the stack addresses
from the main-frame index up to the index obtained by adding $F$ to it.
\end{definition}

\begin{definition}[Pool of encoded indices]
\lean{Light.Compiler.Pool}
\label{Light.Compiler.Pool}
\leanok
\uses{Light.Compiler.cPool}
For each natural number $N$, let $P_N$ be the set of integer codes $c(n)$ for indices
$n$ with $n\le N$.
\end{definition}

\begin{definition}[Constant pool]
\lean{Light.Compiler.Constants}
\label{Light.Compiler.Constants}
\leanok
\uses{Light.Compiler.Pool, Light.Compiler.cNEG, Light.Compiler.cONE, Light.Compiler.cZERO}
For each natural number $N$, let $C_N$ consist of the three distinguished integer codes
$c_1,c_0,c_{-1}$ together with the pool $P_N$ of codes for indices at most $N$.
\end{definition}

\begin{definition}[Scratch locations]
\lean{Light.Compiler.Scratch}
\label{Light.Compiler.Scratch}
\leanok
\uses{Light.Compiler.Temps, Light.Compiler.cD, Light.Compiler.cNFP, Light.Compiler.cRET, Light.Compiler.cRR}
For each natural number $F$, let the scratch set consist of the four distinguished
integer codes $c_D,c_{\mathrm{RET}},c_{\mathrm{RR}},c_{\mathrm{NFP}}$ together with the
temporary codes $T(0,F)$.
\end{definition}

\begin{definition}[Writable locations]
\lean{Light.Compiler.Writable}
\label{Light.Compiler.Writable}
\leanok
\uses{Light.Compiler.Scratch, Light.Compiler.Sizes, Light.Compiler.Stack, Light.Compiler.cFP}
Let $Z$ be a size specification with word width $W$, frame size $F$, dispatcher position
$d$, stack capacity $Q$, and space limit $S$, satisfying its stated fit conditions and
$2(F+1)\le d$. For a natural number $q$, the writable set consists of the scratch
locations for $F$, the frame pointer code, the stack locations at depth $q$ with
capacity $Q$, and all integer addresses $a$ with $0\le a<S$.
\end{definition}

\begin{theorem}[Characterization of temporary addresses]
\lean{Light.Compiler.mem_temps_iff}
\label{Light.Compiler.mem_temps_iff}
\leanok
\uses{Light.Compiler.Temps, Light.Compiler.cADDR, Light.Compiler.cT, Set.mem_ofPred_eq}
\difficulty{4}
For natural numbers $k,F$ and an integer $a$, let $c_T(j)=-(23+4j)$ and let the
temporary address set consist of $-13$ together with the values $c_T(j)$ for $k\le j\le
2F$. Then $a$ belongs to this set if and only if either $a=-13$, or $c_T(2F)\le a\le
c_T(k)$ and $4$ divides $c_T(0)-a$.
\end{theorem}

\begin{theorem}[Characterization of stack addresses]
\lean{Light.Compiler.mem_stack_iff}
\label{Light.Compiler.mem_stack_iff}
\leanok
\uses{Light.Compiler.Stack, Light.Compiler.cStack, Set.mem_ofPred_eq}
\difficulty{3}
For natural numbers $q,Q$ and an integer $a$, write $c(j)=-2j$ for each natural number
$j$, and let the stack address set consist of the values $c(j)$ with $q\le j<Q$. Then
$a$ belongs to this set if and only if $c(Q)<a\le c(q)$ and $2$ divides $a$.
\end{theorem}

\begin{theorem}[Characterization of pool addresses]
\lean{Light.Compiler.mem_pool_iff}
\label{Light.Compiler.mem_pool_iff}
\leanok
\uses{Light.Compiler.Pool, Light.Compiler.cPool, Set.mem_ofPred_eq}
\difficulty{3}
For every natural number $N$ and integer $a$, $a$ belongs to the pool of size $N$ if and
only if it lies between the pool addresses $c_{
m Pool}(N)$ and $c_{
m Pool}(0)$,
inclusive, and $c_{
m Pool}(0)-a$ is divisible by $4$.
\end{theorem}

\begin{theorem}[Temporary addresses lie in scratch space]
\lean{Light.Compiler.temps_subset_scratch}
\label{Light.Compiler.temps_subset_scratch}
\leanok
\uses{Light.Compiler.Scratch, Light.Compiler.Temps}
\difficulty{1}
For natural numbers $k$ and $F$, every temporary address indexed from $k$ through $2F$,
together with the address $c_{\mathrm{ADDR}}$, belongs to the scratch set consisting of
the four designated scratch addresses and all temporary addresses indexed from $0$
through $2F$.
\end{theorem}

\begin{theorem}[Scratch Addresses Are Writable]
\lean{Light.Compiler.scratch_subset_writable}
\label{Light.Compiler.scratch_subset_writable}
\leanok
\uses{Light.Compiler.Scratch, Light.Compiler.Writable}
\difficulty{1}
For any word size $W$, size parameters $Z$, and stack index $q$, every address in the
scratch set for $Z$ belongs to the writable-address set for $Z$ at stack index $q$.
\end{theorem}

\begin{theorem}[Stack cells are writable]
\lean{Light.Compiler.stack_subset_writable}
\label{Light.Compiler.stack_subset_writable}
\leanok
\uses{Light.Compiler.Stack, Light.Compiler.Writable}
\difficulty{1}
For any size configuration $Z$ and stack base $q$, every address in the stack region
from $q$ up to, but not including, the stack capacity of $Z$ is writable at base $q$.
\end{theorem}

\begin{theorem}[Writable sets are antitone in the stack bound]
\lean{Light.Compiler.writable_subset_writable}
\label{Light.Compiler.writable_subset_writable}
\leanok
\uses{Light.Compiler.Writable}
\difficulty{3}
Fix a machine size configuration $Z$. For natural numbers $q$ and $q'$, if $q\le q'$,
then every address writable with stack bound $q'$ is writable with stack bound $q$.
\end{theorem}

\begin{theorem}[Address belongs to writable set]
\lean{Light.Compiler.mem_writable_of_addr}
\label{Light.Compiler.mem_writable_of_addr}
\leanok
\uses{Light.Compiler.Writable, Light.Limits.Addr}
\difficulty{1}
Let $Z$ be a compiler size configuration, let $q$ be a natural number, and let $a$ be an
integer. If $a$ is an address satisfying $0\le a<Z$'s space limit, then $a$ belongs to
the writable address set determined by $Z$ and $q$.
\end{theorem}

\begin{theorem}[Stack address under subtraction]
\lean{Light.Compiler.cStack_sub}
\label{Light.Compiler.cStack_sub}
\leanok
\uses{Light.Compiler.cStack}
\difficulty{3}
For all natural numbers $q$ and $x$, subtracting $2x$ from the stack address with index
$q$ gives the stack address with index $q+x$.
\end{theorem}

\begin{theorem}[Stack address under index addition]
\lean{Light.Compiler.cStack_add}
\label{Light.Compiler.cStack_add}
\leanok
\uses{Light.Compiler.cStack}
\difficulty{3}
For all natural numbers $q$ and $x$, the stack address with index $q+x$, increased by
$2x$, equals the stack address with index $q$.
\end{theorem}

\begin{theorem}[Successive stack addresses]
\lean{Light.Compiler.cStack_succ_add}
\label{Light.Compiler.cStack_succ_add}
\leanok
\uses{Light.Compiler.cStack, Light.Compiler.cStack_add}
\difficulty{1}
For every natural number $q$, the stack address with index $q+1$, increased by $2$,
equals the stack address with index $q$.
\end{theorem}

\begin{theorem}[Stack addresses fit the word range]
\lean{Light.Compiler.Sizes.inRange_stack}
\label{Light.Compiler.Sizes.inRange_stack}
\leanok
\uses{Light.Compiler.Sizes, Light.Compiler.cStack, WordRAM.InRange}
\difficulty{3}
For every word width $W$ and size configuration $Z$, if $j$ is a natural number with
$j\le Z.Q$, then the stack address $-2j$ lies in the signed word range: $-2^W\le
2(-2j)<2^W$.
\end{theorem}

\begin{theorem}[Address values fit the word range]
\lean{Light.Compiler.Sizes.inRange_addr}
\label{Light.Compiler.Sizes.inRange_addr}
\leanok
\uses{Light.Compiler.Sizes, Light.Limits.Addr, WordRAM.InRange}
\difficulty{2}
For every word width $W$, size configuration $Z$, and integer address $a$ with $0\le
a<Z.\mathrm{lim}.\mathrm{space}$, the address lies in the signed word range: $-2^W\le
2a<2^W$.
\end{theorem}

\begin{theorem}[Negative one lies in the word range]
\lean{Light.Compiler.Sizes.inRange_neg_one}
\label{Light.Compiler.Sizes.inRange_neg_one}
\leanok
\uses{Light.Compiler.Sizes, WordRAM.InRange}
\difficulty{3}
For every word width $W$ and every size configuration $Z$ at that word width, $-1$ lies
in the signed word range: $-2^W\le -2<2^W$.
\end{theorem}

\begin{definition}[Machine representation relation]
\lean{Light.Compiler.Rel}
\label{Light.Compiler.Rel}
\leanok
\uses{Light.Compiler.Sizes, Light.Compiler.cFP, Light.Compiler.cNEG, Light.Compiler.cONE, Light.Compiler.cPool, Light.Compiler.cStack, Light.Compiler.cZERO, Light.State, WordRAM.wd}
Let $Z$ be a size specification of word width $W$, let $\sigma$ be a state with local
values $\sigma_{\mathrm{loc}}(x)$ and memory values $\sigma_{\mathrm{mem}}(a)$, let $q$
be a natural number, and let $m$ assign a $W$-bit word to each integer address. The
relation holds when $m$ maps the distinguished codes for $1$, $0$, and $-1$ to those
respective words; maps each pool code for $n\le d$ to the word for $n$; maps the
frame-pointer code to the word for stack position $c_{\mathrm{Stack}}(q)$; maps each
local stack position $c_{\mathrm{Stack}}(q+x)$, for $x<F$, to the word for
$\sigma_{\mathrm{loc}}(x)$; and maps every address $a<S$ to the word for
$\sigma_{\mathrm{mem}}(a)$.
\end{definition}

\begin{theorem}[Constructing a relation from memory agreement]
\lean{Light.Compiler.Rel.of_agreeOutside}
\label{Light.Compiler.Rel.of_agreeOutside}
\leanok
\uses{Light.Compiler.Constants, Light.Compiler.Rel, Light.Compiler.cFP, Light.Compiler.cStack, WordRAM.AgreeOutside, WordRAM.AgreeOutside.cell, WordRAM.wd}
\difficulty{4}
Let $W$ be a word size, let $Z$ specify a layout with $F$ local variables, dispatcher
position $D=Z.\mathrm{disp}$, and space limit, and let $\sigma,\sigma'$ be states with
frame pointers $q,q'$. Suppose $m$ satisfies the storage relation for $Z,\sigma,q$, and
$m'$ agrees with $m$ outside a set $s$ of addresses disjoint from the addresses of the
encoded constants $1,0,-1$ and pool addresses indexed by $n\le D$. Suppose also that
$m'$ stores the stack address for frame pointer $q'$ at the frame-pointer address,
stores $\sigma'_{\mathrm{loc}}(x)$ at stack address $q'+x$ for every $x<F$, and stores
$\sigma'_{\mathrm{mem}}(a)$ at address $a$ for every $a<Z.\mathrm{lim}.\mathrm{space}$.
Then $m'$ satisfies the storage relation for $Z,\sigma',q'$.
\end{theorem}

\begin{theorem}[Relation preserved outside scratch addresses]
\lean{Light.Compiler.Rel.same}
\label{Light.Compiler.Rel.same}
\leanok
\uses{Light.Compiler.Rel, Light.Compiler.Rel.of_agreeOutside, Light.Compiler.Scratch, WordRAM.AgreeOutside, WordRAM.AgreeOutside.read}
\difficulty{2}
Let $W$ be a word size, let $Z$ specify a layout with $F$ local variables per frame,
stack bound $Q$, and space limit, and let $\sigma$ be a state. Suppose $m$ stores the
encoded constants $1,0,-1$ and pool indices $n\le Z.\mathrm{disp}$ at their designated
addresses, stores the stack address for frame pointer $q$, stores
$\sigma_{\mathrm{loc}}(x)$ at stack address $q+x$ for every $x<F$, and stores
$\sigma_{\mathrm{mem}}(a)$ at address $a$ for every $a<Z.\mathrm{lim}.\mathrm{space}$.
If $m'$ agrees with $m$ outside the scratch addresses $-21,-15,-17,-19,-13$, and
$-(23+4j)$ for $0\le j\le 2F$, then $m'$ satisfies the same storage relation for
$Z,\sigma,q$.
\end{theorem}

\begin{theorem}[Preservation of the memory relation]
\lean{Light.Compiler.Rel.beyond}
\label{Light.Compiler.Rel.beyond}
\leanok
\uses{Light.Compiler.Rel, Light.Compiler.Rel.of_agreeOutside, Light.Compiler.Scratch, Light.Compiler.Stack, WordRAM.AgreeOutside, WordRAM.AgreeOutside.read}
\difficulty{2}
Let $W$ be a word size, let $Z$ be a size specification with $F$ local variables per
frame, stack capacity $Q$, and space limit, and let $c$ be a state. Suppose a memory
map $m$ satisfies the relation for $Z$, $c$, and a frame pointer index $q$, where $1
\leq q$. If another memory map $m'$ agrees with $m$ outside the union of the scratch
addresses and the stack addresses indexed by $q+F \leq j<Q$, then $m'$ also satisfies
the relation for $Z$, $c$, and $q$.
\end{theorem}

\begin{theorem}[Updating a local variable]
\lean{Light.Compiler.Rel.setLoc}
\label{Light.Compiler.Rel.setLoc}
\leanok
\uses{Light.Compiler.Rel, Light.Compiler.Rel.of_agreeOutside, Light.Compiler.cStack, WordRAM.AgreeOutside.refl, WordRAM.AgreeOutside.update, WordRAM.wd}
\difficulty{4}
Let $W$ be a word size, let $Z$ specify a layout with $F$ local variables per frame, and
suppose $m$ satisfies the storage relation for a state $\sigma$ and frame pointer $q$.
If $1\le q$, then for any local index $x$ and integer value $v$, the state obtained from
$\sigma$ by replacing its $x$th local value with $v$, together with the memory map
obtained from $m$ by replacing the value at stack address $q+x$ with the encoding of
$v$, satisfies the same storage relation at frame pointer $q$.
\end{theorem}

\begin{theorem}[Preservation of the memory relation]
\lean{Light.Compiler.Rel.setMem}
\label{Light.Compiler.Rel.setMem}
\leanok
\uses{Light.Compiler.Rel, Light.Compiler.Rel.of_agreeOutside, Light.Limits.Addr, WordRAM.AgreeOutside.refl, WordRAM.AgreeOutside.update, WordRAM.wd}
\difficulty{4}
Let $W$ be a word width, $Z$ a size configuration, and let a light state $\sigma$
and a word-RAM memory $m$ satisfy the stated relation at stack position $q$. If $1\le q$
and $a$ is a valid address, meaning $0\le a<Z.\mathrm{lim}.\mathrm{space}$, then for
every integer $v$, updating the light state's memory at the natural-number address
corresponding to $a$ with value $v$ and updating $m$ at address $a$ with the $W$-bit
representation of $v$ again gives related states at position $q$.
\end{theorem}

\begin{theorem}[Reading a bounded memory address]
\lean{Light.Compiler.Rel.read}
\label{Light.Compiler.Rel.read}
\leanok
\uses{Light.Compiler.Rel, Light.Limits.Addr, WordRAM.wd}
\difficulty{3}
Let $W$ be a word size, let $Z$ specify a space limit, and suppose $m$ stores the
encoded value of $\sigma_{\mathrm{mem}}(a)$ at each natural-number address
$a<Z.\mathrm{lim}.\mathrm{space}$. For any integer address $a$ with $0\le
a<Z.\mathrm{lim}.\mathrm{space}$, the value of $m$ at address $a$ is the encoding of
$\sigma_{\mathrm{mem}}(a)$, where the address is viewed as a natural number.
\end{theorem}

\begin{theorem}[First local variable address]
\lean{Light.Compiler.Rel.loc_zero}
\label{Light.Compiler.Rel.loc_zero}
\leanok
\uses{Light.Compiler.Rel, Light.Compiler.cStack, WordRAM.wd}
\difficulty{1}
Let $W$ be a word size, let $Z$ be a layout with $F$ local variables, and suppose $m$
satisfies the storage relation for a state $\sigma$ and frame pointer $q$: in
particular, it stores $\sigma_{\mathrm{loc}}(x)$ at stack address $q+x$ for every $x<F$.
If $F>0$, then $m$ at the stack address $q$ equals the encoding of
$\sigma_{\mathrm{loc}}(0)$.
\end{theorem}

\begin{definition}[Framed machine representation]
\lean{Light.Compiler.Framed}
\label{Light.Compiler.Framed}
\leanok
\uses{Light.Compiler.Rel, Light.Compiler.Sizes, Light.State}
Let $Z$ be a size specification of word width $W$, let $\sigma$ be a state, let $q$ be a
natural number, and let $m$ map integer addresses to $W$-bit words. The representation
is framed when the machine representation relation holds, $1\le q$, and $q+F\le Q$,
where $F$ is the frame size and $Q$ is the stack capacity specified by $Z$.
\end{definition}

\begin{theorem}[Framing is preserved outside scratch addresses]
\lean{Light.Compiler.Framed.same}
\label{Light.Compiler.Framed.same}
\leanok
\uses{Light.Compiler.Framed, Light.Compiler.Rel.same, Light.Compiler.Scratch, WordRAM.AgreeOutside}
\difficulty{1}
Let $W$ be a word size, let $Z$ specify a layout with $F$ local variables per frame and
stack bound $Q$, and let $\sigma$ be a state with local values
$\sigma_{\mathrm{loc}}(x)$ and memory values $\sigma_{\mathrm{mem}}(a)$. Suppose a
memory map $m$ stores the encoded constants $1,0,-1$ at their designated addresses,
stores each pool index $n\le Z.\mathrm{disp}$ at its pool address, stores the stack
address for frame pointer $q$ at the frame-pointer address, and stores
$\sigma_{\mathrm{loc}}(x)$ at stack address $q+x$ for every $x<F$ and
$\sigma_{\mathrm{mem}}(a)$ at address $a$ for every $a<Z.\mathrm{lim}.\mathrm{space}$;
also suppose $1\le q$ and $q+F\le Q$. If $m'$ agrees with $m$ outside the scratch
addresses (the addresses $-21,-15,-17,-19,-13$, and $-(23+4j)$ for $0\le j\le 2F$), then
$m'$ satisfies the same storage and frame-bound conditions for $Z,\sigma,q$.
\end{theorem}

\begin{theorem}[Negative branch jump]
\lean{Light.Compiler.steps_jump}
\label{Light.Compiler.steps_jump}
\leanok
\uses{Light.Compiler.cNEG, WordRAM.CodeAt, WordRAM.InRange, WordRAM.Steps, WordRAM.steps_bltz, WordRAM.wd}
\difficulty{1}
Suppose the code at position $\mathrm{pos}$ begins with a branch instruction that jumps
to $l$ when the tested register is negative. If that register contains the word encoding
$-1$ and $-1$ lies in the signed range $-2^{W-1}\le -1<2^{W-1}$, then one machine step
takes the program counter from $\mathrm{pos}$ to $l$ without changing memory.
\end{theorem}

\begin{theorem}[Compiled assignment to a variable]
\lean{Light.Compiler.steps_setVar}
\label{Light.Compiler.steps_setVar}
\leanok
\uses{Light.Compiler.Framed, Light.Compiler.Rel, Light.Compiler.Rel.same, Light.Compiler.Rel.setLoc, Light.Compiler.Scratch, Light.Compiler.Sizes.inRange_stack, Light.Compiler.Writable, Light.Compiler.cADDR, Light.Compiler.cFP, Light.Compiler.cPool, Light.Compiler.cStack, Light.Compiler.cStack_sub, Light.Compiler.scratch_subset_writable, WordRAM.AgreeOutside, WordRAM.AgreeOutside.mono, WordRAM.AgreeOutside.refl, WordRAM.AgreeOutside.update, WordRAM.CodeAt, WordRAM.CodeAt.tail, WordRAM.Steps, WordRAM.Steps.trans, WordRAM.steps_store, WordRAM.steps_sub, WordRAM.wd}
\difficulty{4}
Let $Z$ be admissible compiler size data of word size $W$, with state $\sigma$
represented in a stack frame at depth $q$, where $1\le q$ and $q+F\le Q$. Suppose $x<F$,
the code at $\mathrm{pos}$ begins with the two instructions that compute the address of
local variable $x$ and store the value from address $c$ there, memory $m$ contains the
word encoding $v$ at address $c$, and $c$ is not the address register. Then after two
machine steps some memory $m'$ represents the state obtained from $\sigma$ by setting
local variable $x$ to $v$, at the same stack depth $q$. The memories agree outside the
writable-address set determined by $Z$ and $q$.
\end{theorem}
